% R008 — Guide to Cultivating Complex Analysis (Bahasa Indonesia)
% Reader-facing translation unit: source/ca-v1.9/ca.tex lines 1324–1353.
% Source release: v1.9 (commit a4d25d9ba37a486a5a020219eb8bd4e6602fd14c).
% Translation provenance: OpenAI Codex gpt-5.6-sol, Ultra, acting on the user's request.
% Preserve source equations, notation, glossary hooks, and footnote structure.

\subsection{Argumen}

Kita mencoba mendefinisikan argumen dari $z = re^{i\theta}$ sebagai
\glsadd{not:arg}%
\begin{equation*}
\arg z
\overset{\text{def?}}{=}
\theta ,
\end{equation*}
tetapi kita menghadapi masalah bahwa jika $\theta$ merupakan salah satu argumen dari $z$,
maka $\theta+2\pi$, $\theta-2\pi$, atau $\theta + k 2\pi$ juga merupakan argumennya
untuk setiap bilangan bulat $k$. Dengan kata lain,
$\arg z$ bukanlah fungsi dalam pengertian klasik, melainkan sebuah
\emph{\myindex{fungsi bernilai banyak}}%
\footnote{Analis yang tidak berkecimpung dalam analisis kompleks terkadang akan menyatakan bahwa
fungsi bernilai banyak tidak masuk akal, tetapi Anda boleh mengabaikan para pengacau itu.}. Definisi
yang tepat adalah
\glsadd{not:arg}%
\begin{equation*}
\arg z
\overset{\text{def}}{=}
\ldots,\theta - 4\pi,
\theta - 2\pi,
\theta ,
\theta + 2\pi,
\theta + 4\pi, \ldots
\end{equation*}
Masih ada satu persoalan kecil. Jika $z=0$, maka $z=0=0 e^{i\theta}$ untuk sembarang
$\theta$. Oleh karena itu, kita hanya mendefinisikan argumen untuk $z$ yang tak nol.
